\section{What the Study Found}
\label{sec:findings}

Twelve developers completed both conditions. This section reports what they did,
what they got right, and what they said---in that order, because the primary
measure (Section~\ref{sec:eval-study}) is errors a participant does not notice,
and everything else is context for interpreting that number. A tool that helps a
developer finish faster while leaving hidden damage is worse than one that slows
them down and catches it, because the developer who submits with confidence has
spent their verification budget and will not look again.

Mozannar et al.\ modelled this budget formally: the dominant cost in AI-assisted
programming is not generation but verification---a developer spends more time
evaluating whether generated code is correct than they saved by not writing
it~\cite{mozannar2024}. The cost scales with the developer's existing
comprehension of the surrounding code, which is exactly the variable our two
conditions manipulate. In the git condition, the developer's comprehension of
which functions serve which request must be constructed from reading diffs. In
the sgt condition, the tool's attribution supplies that comprehension directly,
so the verification cost of a revert preview should be lower---if the attribution
is correct. If it is wrong, the developer inherits a false model and the
verification cost rises above the git baseline, because they must now correct
both the code and their understanding of it. Parasuraman et al.'s model predicts
this asymmetry---automating information analysis (stage~2) without adequate
information acquisition (stage~1) produces exactly the undetected
error~\cite{parasuraman2000}---and the study's primary measure is designed to
catch it: an error that survives the participant's verification is evidence that
the budget was spent on the wrong model.

We pre-registered the measures and the analysis before the first session; the
predictions stated in Table~\ref{tab:measures} were fixed before the data
existed, and the reader can check each one against the result below.

All confidence intervals are 95\% studentized bootstrap over participants
(10{,}000 resamples), because at $n=12$ the sampling distribution of the mean
is not guaranteed to be normal and a $t$-interval would require that assumption.
We report paired estimates (sgt minus git) throughout, so each participant
serves as their own control and individual variation cancels.

\subsection{Participants}

We recruited twelve developers who use a coding agent as part of their ordinary
work. \placeholder{Background: mean years coding, mean years git, git expertise
composite (0--24), agent tools used, AI share of shipped code.} Four groups of
three, fully crossing condition order with project assignment (coursecraft and
confplan). No participant had used \sgt{} before the session.

\subsection{Task outcomes}
\label{sec:findings-outcomes}

Figure~\ref{fig:outcomes} shows the three primary outcome measures, paired by
participant. The measure the paper committed to first is collateral damage:
tests broken outside the set the participant was asked to remove.

\begin{figure*}[t]
  \centering
  % \input{figures/fig-outcomes}   % PairedEstimation chart, 3 panels
  \placeholder{Figure: PairedEstimation with panels R1 provenance (0--3),
  R2+R3 removal (0--4 rubric points), and collateral damage (inverted axis).
  Lines connect each participant's git and sgt scores. Thick bars are condition
  means. Below each panel, the paired mean difference with its bootstrap
  distribution and 95\% interval.}
  \caption{Task outcomes by participant and condition. Every participant appears
  twice, joined by a line. Below each panel, the paired mean difference (sgt
  minus git) with its bootstrap distribution and 95\% studentized interval.
  Collateral damage is plotted with the axis inverted, so up is better in all
  three panels.}
  \Description{Three paired estimation panels showing per-participant scores on
  R1 provenance, R2+R3 removal, and collateral damage, each with a bootstrap
  difference estimate below.}
  \label{fig:outcomes}
\end{figure*}

\paragraph{Provenance (request~1).}
Three closed questions, scored against a key.
\placeholder{Git mean, sgt mean, paired difference, CI, $d_z$, n.}
The confidence calibration (stated confidence minus proportion correct)
was \placeholder{git calibration, sgt calibration, difference} --- positive is
overconfident. \placeholder{Interpret: did sgt participants know more, or merely
believe they knew more? If calibration is flat while score rises, the tool
helped them be right. If calibration also shifts, they were also more sure of
what they got right --- or more sure of what they got wrong.}

\paragraph{Reach prediction (Foresee block).}
Two trials, each scored as F1 against measured ground truth (call-graph reach).
\placeholder{Per-trial blind, checked, gain means by condition. Paired gain
difference (sgt minus git), CI. Confidence calibration at both stages.}

\noindent Trial F1 (agenda export, narrow reach): \placeholder{blind, checked,
gain per condition.}
Trial F2 (slot-comparison normalisation, broad reach): \placeholder{blind,
checked, gain per condition.}

\placeholder{Interpret: (1) Is gain positive under sgt?  If so, the
representation moved the participant's prediction toward truth. (2) Is gain
larger under sgt than under git?  That is the between-condition test. (3) Did
any participant show high blind confidence paired with low blind accuracy under
sgt?  That is the inferred-label failure mode---confidently wrong rather than
helpfully wrong---and if it appears on F2 (the pre-registered risk case), it is
reported as a U1 failure.  (4) Does the narrow-reach trial (F1) produce
different gain than the broad-reach trial (F2)?  If gain is asymmetric, the
representation helps in one direction (pulling predictions down or pushing them
out) but not the other, which says something specific about which kind of
provenance question the tool serves.}

\paragraph{Removal and correction (requests~2 and~3).}
Task score (rubric points): \placeholder{git mean, sgt mean, difference, CI}.
Collateral damage (failing tests outside the target):
\placeholder{git mean, sgt mean, difference, CI}.
\placeholder{Interpret: the critical question is whether the sgt condition's
consequence preview let participants avoid the collateral that the git condition
could not see coming. If collateral is equal, recovery is something developers
already do reliably and Section~\ref{sec:scenarios} overstates the difficulty.}

\paragraph{Time and cap pressure.}
\placeholder{Proportion hitting the cap per condition. Mean active time.
Manipulation check: felt time pressure response distribution per condition.}
We report time because a reader would want it, and we would not publish a time
difference alone: a tool that is correct and unbearably slow is not adopted, but
a tool that is fast and wrong is worse.

\subsection{How the work was done}
\label{sec:findings-process}

Outcome measures say whether a tool helped; process measures say how. Two
conditions can produce the same error rate through different strategies---one
because the developer checked carefully, the other because the tool caught
things the developer missed---and the distinction matters for design. The
process telemetry also grounds the preference data: a developer who says ``I
trusted it'' while the telemetry shows frequent verification is better
calibrated than one whose telemetry confirms the trust.

\begin{figure*}[t]
  \centering
  % \input{figures/fig-process}   % ProcessSignature chart
  \placeholder{Figure: ProcessSignature with (a) share of action categories
  across normalized request time, pooled per condition, with one strip per
  participant beneath; (b) action bigrams that most distinguish the conditions,
  by weighted log-odds with informative Dirichlet prior.}
  \caption{(a)~Share of action categories across normalised request time, pooled
  per condition, with one strip per participant beneath.
  (b)~Action bigrams that most distinguish the conditions (weighted log-odds,
  informative Dirichlet prior; positive leans sgt).}
  \Description{A process signature chart showing action category distributions
  over time for both conditions, plus distinguishing bigram sequences.}
  \label{fig:process}
\end{figure*}

The process telemetry recorded every command, prompt, and edit action,
categorised into nine kinds (orient, inspect, search, prompt, agent edit, manual
edit, history operation, verify, recover). The first three are navigation (the
developer is looking for something); the next three are transformation (the
developer or assistant is changing something); the last three are quality
assurance (the developer is checking or undoing). Figure~\ref{fig:process}a
shows how those three phases distribute over normalised request time.

\paragraph{Verification ratio.}
\placeholder{Git mean, sgt mean, difference, CI.}
The verification ratio is the proportion of actions in the ``verify'' category
(running tests, reading diffs, inspecting output) relative to all actions within
a request. The direction of this measure is ambiguous by design: a higher ratio
in sgt could mean the preview gave participants something concrete to check
(supporting appropriate reliance), or it could mean the tool's unfamiliarity
forced more checking (friction). A lower ratio in sgt could mean participants
trusted the consequence display and skipped verification (efficient trust, or
dangerous complacency). The self-report item q14 (accepted unreviewed work,
reversed; Section~\ref{sec:findings-perception})
disambiguates: if verification drops AND q14 rises, participants skipped checks
they felt they should have made.
\placeholder{Report the numbers and which disambiguation holds.}

\paragraph{Prompt specificity.}
\placeholder{Git mean (0--3), sgt mean, difference, CI.}
Each prompt to the AI assistant was coded on a four-point specificity scale: 0
(no target named) through 3 (names a file, a function, and what to do). The
coding was applied to the raw prompt text by the taxonomy of
Section~\ref{sec:eval-study}, not by a human judge, so inter-rater reliability
is not a concern and every prompt across both conditions received the same
treatment. The measure tests whether the tool's vocabulary---feature names,
checkpoint labels, symbol identifiers visible in the map---entered the
developer's prompts. Two opposing mechanisms are plausible. If sgt's names gave
participants concrete referents they otherwise lacked, specificity should rise:
a developer who can see that the waitlist is feature \texttt{f-7ab21c}
writes ``revert f-7ab21c'' (level~3) instead of ``undo the waitlist stuff''
(level~1). If the tool substituted for the developer's own precision---they
delegated navigation to sgt and therefore never acquired the vocabulary
themselves---specificity should fall, which is the scaffolding-removal concern
Kazemitabaar et al.\ measured in learners~\cite{kazemitabaar2025}. The
distinction matters for the design's long-term effect: a tool that raises
specificity is augmenting the developer's coordination with the assistant; a
tool that lowers it is creating a dependency the developer cannot sustain when
the tool is absent.
\placeholder{Report which direction holds and connect to the vocabulary theme
of Section~\ref{sec:findings-reflection}.}

\paragraph{Wrong turns.}
\placeholder{Git sum, sgt sum, per-condition.} A wrong turn is a history
operation followed by a recovery action (undo, restore, or manual repair) within
two minutes---the participant tried something, saw it was wrong, and backed out.
The direction of this measure is the sharpest test of whether the pilot's
mutation failures persist in the fixed version. If wrong turns are higher in
sgt, participants reached the write layer and found it still unreliable: the
tool stopped them (good) but stopping is work, and enough stops constitute a
wrong-turn tax that erodes the time the comprehension layer saved. If wrong
turns are lower in sgt, participants either avoided the write layer entirely
(using sgt only to read, then acting in git) or the fixes actually stabilised
it. The telemetry distinguishes these two explanations: an absent write layer
shows up as zero history\_op events of the mutating kind, while a stable write
layer shows up as history\_op events without recovery following them.
\placeholder{Report numbers and which explanation holds.}
In Ferino et al.'s reliance-control framework~\cite{ferino2026}, a wrong turn
followed by recovery is a developer exercising control at the appropriate
moment---they detected a problem and corrected it. A wrong turn followed by
abandonment (switching to git for the remainder) is control escalation: the
developer decided the tool's level of reliability does not merit continued
engagement. The telemetry distinguishes these, and the second pattern is the
one that would count against the design, because a tool that drives developers
away from its write layer on first contact has failed to build the reliance its
read layer earned.

\paragraph{Distinguishing sequences.}
\placeholder{Top 3--5 bigrams that lean sgt, top 3--5 that lean git, with
log-odds and count. Interpret each: what does the sequence mean as behaviour?
For example, if ``history\_op → verify'' is strongly sgt, participants checked
what the tool did. If ``search → search'' is strongly git, participants were
hunting through files.}

The bigram comparison uses weighted log-odds with an informative Dirichlet prior
(Monroe et al.~\cite{monroe2008}), which controls for frequency: a rare bigram
that appears only in one condition produces a high ratio but a low weight,
preventing one participant's idiosyncratic strategy from dominating the
comparison. A positive $z$-score leans sgt; negative leans git. We report
bigrams above a minimum corpus frequency of five, and state how many rare
sequences were dropped, because silent truncation reads as ``we looked at
everything'' when we did not.

In information foraging terms~\cite{pirolli1999}, the two conditions differ in
how rich each information ``patch'' is. Feature names and consequence previews
are scent: proximal cues that predict what an operation will reveal before the
developer commits to it. The n-gram analysis tests this directly. If the sgt
condition shows fewer search$\to$search cycles and more
history\_op$\to$verify sequences, the tool enriched the history record as a
foraging patch---developers spend less time in between-patch navigation and
more in within-patch exploitation, which is the 35\% of maintenance time Ko
et al.\ measured as pure navigation cost~\cite{ko2006}.

Two further patterns would be informative. First,
inspect$\to$history\_op in sgt would mean that seeing the record's structure
gives participants enough confidence to act without further search---the
attribution serves as Pirolli and Card's ``information scent,'' a proximal cue
of distal content. Second, history\_op$\to$recover in sgt would mean write
operations failed often enough that participants had to undo them---the mutation
layer's residual instability showing up as a behavioural signature. If both
conditions show the same cycling, feature names carry no more scent than commit
messages and the vocabulary is not reaching the developer's search behaviour.

Letovsky's comprehension model offers a finer prediction~\cite{letovsky1987}.
Programmers reading unfamiliar code generate questions (``why is this here?'',
``what does this call?'') and conjectures (``I think this handles retries''),
and the questions drive their navigation. In the git condition, the
provenance question---\emph{which request produced this function?}---has no
answer in the record, so the developer must generate it as a conjecture, test
it against the code, and revise. In the sgt condition, the attribution answers
it directly: the developer reads the feature label, confirms or renames, and
moves to the next question. The observable difference is a shift from \emph{generative} inquiry
(conjectures, with their risk of wrong turns) to \emph{verificatory} inquiry
(checking a provided answer), which shows up as fewer search events and a
different sequence structure. The telemetry can distinguish these: generative
inquiry shows up as search$\to$search$\to$inspect sequences where the
developer builds understanding incrementally; verificatory inquiry shows up as
inspect$\to$history\_op, where a single inspection provides enough confidence
to act. If the sgt condition shifts the balance toward verification while
maintaining accuracy, the tool's attribution is doing what Letovsky's model says
a programmer's own memory does when reading code they wrote themselves: supplying
answers to the questions before the search has to generate them.

\subsection{What the two conditions felt like}
\label{sec:findings-perception}

\begin{figure*}[t]
  \centering
  % \input{figures/fig-perception}  % LikertDiverging chart
  \placeholder{Figure: LikertDiverging showing 14 HLAC items grouped in 4
  blocks (Finding your way around, Changing things, What you came away with,
  Working with the assistant), one panel per condition, with paired mean
  differences and CIs as dots and bars.}
  \caption{Perceived experience on fourteen 7-point items (HLAC battery),
  grouped into four ad-hoc blocks. Reverse-coded items (marked~$\hookleftarrow$)
  are recoded so agreement always means better. Dots are paired mean differences
  (sgt minus git); bars are 95\% bootstrap intervals.}
  \Description{A diverging Likert chart with 14 items across two conditions,
  showing response distributions and paired mean differences.}
  \label{fig:perception}
\end{figure*}

The HLAC battery (History Legibility and Agent Collaboration) asked fourteen
7-point items after each half, grouped into four blocks.
No validated scale exists for history comprehension in AI-assisted development;
general trust-in-automation instruments~\cite{jian2000} measure disposition
toward automated systems but cannot distinguish provenance legibility from
consequence awareness from retained mental model, which are the three mechanisms
this paper claims.  We therefore wrote items that map one-to-one to specific
design claims, report them as ad-hoc Likert-type items rather than construct
scores, and ground interpretation in the exit interview where participants state
what they experienced in their own terms.
Figure~\ref{fig:perception} shows the full response distribution
and the paired difference per item.

\paragraph{Finding your way around (C1).}
Four items: found when it changed, found why it changed, saw the whole piece of
work, guessed at names (reversed).
\placeholder{Item-level means and paired differences. Interpret: these map
directly to claim C1. If sgt participants found provenance more legible, these
items should shift. The reversed item (guessing at names) tests whether inferred
labels helped or confused.}

\paragraph{Changing things (C2).}
Four items: knew what else it would touch, removed a feature safely, recovered
from mistakes, surprised by the result (reversed). These four map to the
consequence-preview mechanism of Section~\ref{sec:design}: the tool shows,
before applying a revert, which other functions still depend on something being
removed. That display is the bridge between comprehension (seeing what would
break) and action (deciding to proceed). If participants in the sgt condition
rated ``knew what else it would touch'' higher, the preview reached them---they
read the consequence list and it calibrated their expectation. The reversed
item (q12, ``surprised by the result'') tests whether that calibration held
through execution: a participant can report knowing what would happen and still
be surprised if the tool's execution differed from its preview.
\placeholder{Item-level means and paired differences. If collateral is lower
AND q4 (knew consequences) shifts, participants both did better and knew they
would---the preview set accurate expectations AND the execution honoured them.
If q4 shifts but collateral is flat, participants felt more informed without
producing fewer errors, which is Bansal et al.'s finding restated: explanations
increase confidence without increasing accuracy~\cite{bansal2021}.}

\paragraph{What you came away with (C3).}
Two items: clear picture of the project (q7), would get back up to speed (q13).
This is the weakest block in the battery and we state why rather than papering
over it. Two items cannot establish a construct. They are here because the claim
they probe---that the tool leaves the developer with a transferable mental model
of the codebase---is central to the paper's argument but too diffuse to measure
with four items in a two-hour session. A developer's ``picture of the project''
accumulates across days, and what two hours can measure is at best whether the
tool planted the seed of one. The study design compensates by grounding whatever
signal these items produce in the exit interview, where a participant can say what
they came away with in their own terms rather than rating an abstraction.
\placeholder{Item-level means and paired differences. If both items shift toward
sgt, the tool's attribution gave participants more than task-local comprehension---it
left them with a model of the project's structure they believe they could return
to. If neither shifts, the twenty-minute window was too short for the seed to take
root, and a longitudinal deployment study would be needed to test C3. Report
honestly either way: a null on two items in twenty minutes does not refute the claim,
but neither does it support it, and the reader should know that.}

\paragraph{Working with the assistant (Q4).}
Four items: directed the assistant precisely, checked the assistant's work,
accepted unreviewed work (reversed), fought the tool (reversed). The two
reversed items serve as honesty valves: a participant who straight-lines
``agree'' on every positive item is caught by the negative ones. But they also
measure real phenomena. ``Accepted unreviewed work'' (q14) is the
over-reliance signal: a developer who trusts the tool's consequence preview
and skips independent verification has delegated the check to the tool, and the
tool's guarantee is only as strong as its rebuild fidelity (Section~\ref{sec:eval-exact}).
If sgt wins on comprehension items (C1, C2) while q14 also rises in sgt, the
tool improved understanding \emph{and} reduced the verification that keeps
understanding honest---which is Bansal et al.'s finding that explanations
increase acceptance regardless of correctness~\cite{bansal2021}.
``Fought the tool'' (q10) is the interface-friction signal. The pilot's
vocabulary-mismatch theme predicted that even a fixed tool would impose a
learning cost---the concepts are unfamiliar, the operations are new, and twenty
minutes is not enough to build fluency. If q10 is higher in sgt while outcome
measures also favour sgt, the tool is useful but abrasive, and a developer who
adopted it would pay a front-loaded cost the study's two-hour window cannot
observe amortising.
\placeholder{Report item-level means and paired differences. State which pattern holds.}

\paragraph{Workload and usability.}
NASA-TLX (raw, unweighted): \placeholder{git mean, sgt mean, difference, CI.}
UMUX-Lite (0--100): \placeholder{git mean, sgt mean, difference, CI.}
TLX measures what the task cost; UMUX-Lite measures whether the developer
thought the setup met their requirements. The two instruments answer different
questions and can point in opposite directions. A tool that is useful but
demanding (higher UMUX-Lite, equal or higher TLX) is a tool worth learning: it
delivers capability that justifies its cognitive cost, and the cost may
amortise with fluency. A tool that is easy but insufficient (lower TLX, lower
UMUX-Lite) failed to provide what the developer needed, regardless of how
little it asked. The combination the pilot would predict---based on the
vocabulary-mismatch finding and the mutation failures---is higher TLX in sgt
(the unfamiliar interface costs effort) with higher UMUX-Lite in sgt (the
consequence preview and attribution answer questions git cannot). If TLX is
flat while UMUX-Lite rises, the tool did not impose measurable additional
load---which would mean the learning cost the pilot predicted did not
materialise in a twenty-minute window, either because the fixes removed the
friction or because the task was short enough to stay below the threshold where
unfamiliarity binds.
\placeholder{Report numbers and which pattern holds.}

\subsection{Preference}
\label{sec:findings-preference}

After both halves, each participant compared the two setups on seven job-level
items (five-point scale, symmetric, neither tool named) and two hypothetical
scenarios. The seven items span a gradient from low-stakes contexts (a throwaway
prototype) to high-stakes ones (the participant's own production codebase), so a
participant who differentiates across contexts---preferring sgt for their real
code but indifferent on a throwaway---is telling us that the tool's value scales
with the cost of getting things wrong, which is the claim Section~\ref{sec:scenarios}
makes. A participant who straight-lines the same answer across all seven items is
either genuinely undifferentiated (both tools felt the same) or not engaging with
the gradient, and the two hypothetical scenarios (``a codebase you inherited from
a colleague'' and ``a fresh project you are starting alone'') separate the two:
the inherited case is the paper's scenario and the solo case is the control where
the developer holds the decomposition themselves.
\placeholder{Report per item: distribution, recoded mean. Do participants
differentiate across the gradient? If so, where does the preference cross zero---at
what level of stakes does the tool's cost (learning, friction) start being
justified by its value (comprehension, fewer errors)?}

The ``what made the difference'' multi-select gives the reasons in the
participants' own language. The ``what would put you off'' multi-select collects
the price in the same terms. These two together answer a question no Likert item
can: whether the developer's stated reason for preferring or rejecting the tool
matches the mechanism the design predicted. If the reasons cluster around
comprehension (``I could see what belonged together,'' ``the consequence preview
helped'') the read layer is carrying the value. If they cluster around mutation
(``I could take things out cleanly'') the write layer is carrying it despite the
pilot's predictions. If the costs cluster around trust (``not sure what it had
done'') the pilot's strongest finding survived the fixes.
\placeholder{Report both distributions. Map each selected reason back to the
design commitment it reflects.}

The preference data answers a question the outcome data cannot: whether a tool
that helps is also a tool a developer would install. Lee and See's framework
predicts that the two can diverge---a tool whose occasional failures are
dramatic enough to remember can be distrusted even when its aggregate
performance is positive, because the availability heuristic weights vivid
single events over base rates~\cite{lee2004}. Two items in the cost
multi-select test this directly: ``not being sure what it had done'' (trust
erosion surviving from the pilot's strongest finding) and ``it grouped or named
things in ways I disagreed with'' (commitment~4, inferred names, costing more
than the pilots predicted). If either appears at high frequency while the
outcome measures favour sgt, the tool is useful and unwanted---which is
Bu\c{c}inca et al.'s finding for the most effective cognitive forcing
functions, restated as a product problem~\cite{bucinca2021}.

\paragraph{Exit interview: thematic analysis.}
After the preference items, each participant answered three open questions: what
they noticed about how the two setups differed, what they would want changed in
the setup they preferred less, and whether they thought about the code
differently after the session than before it. The analysis follows Braun and
Clarke's reflexive thematic analysis~\cite{braun2006}: transcripts are coded
inductively, codes are grouped into themes, and themes are checked against the
design commitments to see whether the participant's account of their experience
maps onto the mechanism the design predicted or onto a mechanism the design did
not anticipate.

Four a priori themes derive from the four commitments: (1)~\emph{granularity
mismatch}, where the participant describes the tool's units as too fine, too
coarse, or misaligned with what they wanted to name; (2)~\emph{consequence
awareness}, where the participant describes knowing (or failing to know) what
an operation would affect before it ran; (3)~\emph{trust trajectory}, where the
participant describes a moment that changed how much they believed the tool's
reports; and (4)~\emph{vocabulary adoption}, where the participant uses the
tool's names (feature labels, checkpoint names, symbol identifiers) in their own
description of the codebase rather than speaking in file or function terms. A
fifth theme is reserved for mechanisms the design did not predict: any coherent
account of value or cost that does not map onto the four commitments is coded
separately and reported as a finding about the design space rather than a finding
about this system.
\placeholder{Report: which themes appeared, how many participants generated each,
and the one or two quotations that most sharply illustrate the mechanism. Map
each theme back to its design commitment. If theme~5 appears with more than two
participants, it is the finding that matters most, because it means the
developer's experience of the tool is shaped by something the designer did not
control.}

\subsection{Reflecting on the design commitments}
\label{sec:findings-reflection}

The pilots predicted a specific shape: comprehension holds, mutation fails
(Section~\ref{sec:pilot-lessons}). Twenty-one fixes and four design corrections
separated the pilot version from the version participants used. The question this
subsection answers is whether those fixes changed the shape or merely moved it.

\paragraph{Did the read/write split survive?}
The pilot participant's exit verdict was ``for reading history, yes, starting
today; for anything that writes, no.'' The study's R1 (provenance) is a pure
read task and R2+R3 (removal/correction) require writing. A split that
survives twelve participants is a finding about the design rather than about
one person's session: it says that the representation (edits filed under
requests) is independently valuable even when the operations over it (revert,
restore) are not yet trustworthy, because the representation answers questions
git cannot express---\emph{which request produced this function, and what else
came with it}---without requiring the developer to act on the answer through
any particular tool. A split that closes (R2+R3 also improved) is a different
finding: the pilot's mutation failures were implementation defects rather than
design costs, and fixing them unlocked the operations the representation makes
possible.
\placeholder{If R1 shows a clear advantage and R2+R3 does not, the split
survived the fixes. If R2+R3 also improved, the fixes changed the story ---
the mutation layer's defects were the bottleneck, not the design.
If neither improved, the pilot's read benefit was an artefact of one participant
rather than a property of the representation.}

\paragraph{Did trust recover?}
The pilot's strongest finding was that a single false positive (green checkmark
over corruption) eroded trust for the remainder of the session, and the trust
did not recover within that session even after correct operations followed. Lee
and See's framework explains why: trust calibration requires multiple consistent
observations to revise, and a single session provides too few trials after a
failure for recalibration to occur~\cite{lee2004}. The version participants used
fixed all twelve defects the pilot found, so the question shifts from ``can
trust recover within a session'' to ``does a tool that never produces the
failure build trust from scratch in the time available.'' The HLAC item
``surprised by the result'' (q12, reversed) and the preference item ``not being
sure what it had done'' are the signals. These two measure different things: q12
asks whether the tool's behaviour matched expectations (a property of the
current session), while the cost item asks whether the developer would carry
uncertainty forward (a property of anticipated future sessions). A tool can
score well on q12 (no surprises today) while still triggering the cost item
(I would not trust it tomorrow), and that combination would mean the pilot's
trust damage left a residue in the design's reputation that the fixes did not
fully erase.
\placeholder{If q12 is flat, participants in the fixed version were not
surprised. If ``not being sure'' still appears frequently in the cost
multi-select, trust did not fully recover despite the fixes --- meaning the
tool's legibility gap is structural (commitment~4, inferred names) rather than
a bug in the mutation layer.}

\paragraph{Did the vocabulary reach the developer?}
The pilot's vocabulary-mismatch theme (four counting units, selector namespaces
varying per verb) predicted that even the comprehension layer would be hard to
reach. The study version unified to one counting unit (symbols) and one
selector namespace. The HLAC item ``guessed at names'' (q11, reversed) and the
process measure ``prompt specificity'' are the signals.

This question matters beyond usability. A version control system that infers
feature names is proposing a shared vocabulary for a codebase---the same
function Furnas et al.\ showed designers cannot predict for their
users~\cite{furnas1987}. If a developer adopts the tool's names into their
prompts, those names become the vocabulary through which they and the assistant
coordinate about the code. The tool has then changed how the developer talks about what they are
doing, which is a stronger form of value than a faster workflow.

\placeholder{If q11 improved and prompt specificity is higher in sgt, the
tool's vocabulary entered the developer's prompts --- they used feature names to
direct the assistant. If specificity is lower, the tool substituted for the
developer's precision rather than augmenting it, which is the over-reliance
concern of Bansal et al.~\cite{bansal2021}.}

The vocabulary question also tests a specific design bet about the
over-decomposition. The tool presents four to six structural communities where
the developer expects one named request, and the design bets that developers will
adopt the finer vocabulary for the same reason they adopt file names: a
structural boundary that stays stable across sessions is more useful as a
referent than a purpose boundary that shifts with rephrasing. If the vocabulary
enters prompts, developers found the subordinate-level labels (``Rate Limiting,''
``Price Cache,'' ``Row Accessors'') useful as handles for directing the
assistant, even though none of them names the request they typed. If it does
not, the over-decomposition is finer than the developer wants to
\emph{talk about}, and a vocabulary nobody adopts fails Clark's grounding
test~\cite{clark1996}, regardless of whether it is structurally
correct. The pilot's git participant provides the sharpest test: they
asked for ``a question about a symbol, not a file,'' which is the structural
level, not the request level. If study participants in the sgt condition adopt
feature names at comparable frequency, the design bet is confirmed: developers
prefer structural referents when they are available, and the mismatch between
tool grain and intent grain matters less than whether the tool's vocabulary is
learnable.

\paragraph{What the preparation found.}
Building the study repositories and running three pilots produced thirty-five
defects, and the pattern across them is informative independent of their
individual fixes. Twenty-two entered through the write path (revert, restore,
save, fulfill) and zero through the read path (log, show, map, find). Twelve
of the twenty-two shared a structural cause: a mutation that removes or
rearranges records leaves metadata (layout chains, anchor positions, residue
sentinels) pointing at records that are no longer present, and the rebuild
produces invalid output from valid arithmetic. This is commitment~2 (set-based
versions) failing at the seam it was designed to hide: files are rebuilt from
sets of edits, and the rebuild is exact only when every piece of metadata that
governs placement is consistent with the set it reads. The remaining ten write failures shared a different cause---one worth naming
as a class: the silent success. A command reports that it finished, or a view
displays a confident summary, and neither did what it claimed. A revert printed ``removes 0 edits'' and
then rewrote two files. A feature named \texttt{Waitlist Queue} held the CLI
scaffold and no waitlist code, sitting one row above the real waitlist. A
save recorded 366 of 370 operations with no provenance because the clustering
signal read a field that the save path never populated---the quality degraded
specifically on the workflow the study asks participants to use.
Ten instances of this pattern surfaced across the preparation, all
recoverable (via \cmd{undo} in every case), all telling the developer something false about what to do next.
The common mechanism is that a tool whose entire pitch is legible consequences
can produce illegible consequences from correct internal arithmetic: every
number was computed honestly, and the sentence it was placed into was wrong.

The silent-success class is architecturally distinct from a crash. A crash
teaches the developer something true about the system's limits: this operation
failed, and the developer can reason about which operations share its structure
and might fail the same way. A silent success teaches something false: this
operation worked, which the developer generalises to other operations of the
same kind. Tang et al.\ found that across 20{,}574 real-world coding-agent
sessions, inaccurate self-reporting---the agent claiming success on work it did
not complete correctly---accounts for 22.6\% of all developer--agent
misalignment episodes, and its share is \emph{rising} even as other failure
classes decline~\cite{tang2026agentfailures}. A tool that adds its own
self-reporting failures to a context already saturated with them is compounding
the ecosystem's dominant residual failure mode rather than compensating for it.
Shen and Tamkin found that the skill most degraded by AI assistance
is debugging---the ability to recognise when code is wrong and
why~\cite{shen2026skills}---and a tool that undermines that same skill by
confirming a wrong state as correct is compounding the very deficit it was
built to address. A version control system that exists because the developer
can no longer hold a model of what the agent wrote must not itself produce
states the developer cannot distinguish from correct ones, because the
developer's degraded capacity to detect wrongness is the tool's entire
reason for existing. The gap between the twelve metadata failures (which
produce visibly wrong output: code that does not parse) and the ten
silent-success failures (which produce invisibly wrong output: code that
parses but does not do what the report claimed) is therefore not a severity
spectrum but a category boundary: the first class fails safely, the second
fails in the specific way the tool was designed to prevent.

If the study's process telemetry shows history\_op events followed by recovery
in the sgt condition, these residual write-path failures are the likely
mechanism. If it shows history\_op events followed by verify (and no recovery),
the twenty-one fixes closed the path and participants reached the stable
operations that survived the preparation intact.

\paragraph{What the robustness harness measured.}
The harness ran 289 randomised sequences totalling 10{,}237 operations across
nineteen repository shapes (fourteen synthetic fixtures, five open-source
repositories \sgt{} had never recorded). Each sequence applies a random
interleaving of saves, reverts, restores, undos, merges, and renames, then
checks every file against the byte-identical reference git holds. A sequence
passes only if every operation leaves every file matching; a single differing
byte on a single file is a violation.

Of 289 sequences, 145 (50\%) contained at least one violation on a
user-issuable operation---meaning a developer working through the same sequence
would have encountered a file whose contents silently differed from what the
tool's report described. At the per-operation level the rate is lower: 340 of
10{,}237 (3.3\%), concentrated on entity-target operations (8.1\% of 2{,}447)
and layout-target operations (4.1\% of 2{,}308), with non-entity, non-layout
operations at 0.8\% (46 of 5{,}482). No sequence produced a traceback---the
system never crashed---and no sequence lost data (the append-only store held
throughout). What it produced instead was the failure mode the preparation
identified: operations that succeed while leaving the files in a state the
report does not describe.

The 50\% per-session rate is the number a developer meets. A developer who uses
\sgt{} for a morning of interleaved saves and reverts has a coinflip chance
of producing a file that silently disagrees with the tool's last report. The
3.3\% per-operation rate describes a different experience: most commands work,
and the developer has no way of predicting which one will not without running
the verification the tool was built to replace. The gap between the two numbers
is the gap between ``it usually works'' and ``I cannot trust any single result
without checking,'' and the second is what governs adoption. Four of the five
real (non-fixture) repositories produced at least one violation, so the rate
is not an artefact of adversarial synthetic inputs.

\paragraph{What the transcript evaluation measured.}
Coding pairwise agreement between \sgt{}'s grouping and four agent transcripts
found that the dominant disagreement was not a clustering error but a
ground-truth problem. In CodeNav (30 coded pairs), 15 were ``consecutive turns
of one feature''---the transcript's turn boundary fell inside a continuous piece
of work, so neither system was wrong; they disagreed about where a boundary
should be. In eico (30 coded pairs), 15 were ``continuation tokens, not
intents'': the transcript carried requests like ``resume,'' ``move on,'' and
``ok sure,'' whose footprint is everything an autonomous agent then did over a
long uninterrupted run. A pairwise F1 computed on a history whose ground truth
is 50\% continuation tokens measures the grain of the ground truth rather than
the quality of the clustering. The share of contentless turns varies from 0\%
to 88\% across the four repositories, and any aggregate metric that pools them
without reporting the contentless fraction hides a confound larger than the
signal.

The one mechanism that \emph{is} an error in the clustering---intent sharded
across leaves (coded F32)---appeared in both repositories and has a consistent
shape: two leaves carry the same label and the same intent, separated only
because their symbols live in different directories, with no interior node
uniting them. In eico, F32 accounts for 11 of 30 coded disagreements. The
fix is structural: a merge pass that detects identically-labelled sibling
leaves would eliminate it without changing any clustering parameter. That it
was not implemented before the evaluation is evidence that the evaluation
found a defect the developer's own use did not surface---precisely the role
the transcript corpus was designed to play.

The thematic pattern across both repositories reflects back on a design choice
Section~\ref{sec:names} defends. The tool's grouping disagrees with the ground
truth for two structurally different reasons: the ground truth is at the wrong
grain (30 of 60 coded pairs across both repositories), or the tool shards one
intent across leaves that share a label but not a directory (14 of 60). The
first is not fixable by improving the clustering---it is a property of the
ground truth, which counts ``resume'' as a boundary---and the second is fixable
by a single post-processing pass. No coded disagreement traced to the coupling
signal itself producing a wrong boundary between genuinely different intents.
The implication is that commitment~4 (inferred names from structural coupling)
holds as a grouping mechanism, but the grain it recovers is finer than what a
developer would call ``one thing I asked for,'' which is the over-decomposition
the conclusion reports. A tool builder choosing between coupling-based grouping
and request-boundary grouping should know that coupling gives stable structural
communities whose boundaries reflect code architecture, while request boundaries
give unstable units whose meaning depends on whether the developer typed
``implement caching'' or ``resume.''

Symbol coverage across the four repositories ranged from 49\% (semipy-package,
83 ground-truth symbols, 41 known) to 79\% (eico, 357 symbols, 282 known).
The gap represents symbols the tool literally cannot answer questions about:
functions that exist in the ground truth but appear in no \sgt{} record,
because the commit that introduced them predates the mining or because the
miner's tier filter excluded their file. A developer who asks ``what produced
this function?'' about a symbol in the gap receives silence rather than a wrong
answer, which is the correct failure mode (honest ignorance rather than
confident error), but it means the tool's coverage is a ceiling on the
fraction of provenance questions it can answer rather than an approximation.

\paragraph{What the evaluation taught about measuring itself.}
The evaluation journal's most consequential finding was not about the tool but
about the measurement. On August~15, a derivability check that had passed eight
times was re-run after a fix, and the result improved so sharply that the
evaluator stopped trusting it. Comparing the trees revealed that the fixture
shipped at signals version~3 while the census copies in \texttt{/tmp} had been
rebuilt at version~2---34 leaves where the fixture had 21, every feature label
the check sequences typed existing only in the rebuilt copy, and not one of the
twelve recorded command sequences capable of running against the shipped
artefact. The census, every derivability check, and the stop-gate result all
described an artefact that existed in \texttt{/tmp} and nowhere a participant
would encounter.

The failure has a specific mechanism: the analysis tool (\texttt{sgt refresh})
quietly rebuilds the store when it detects a version mismatch, the evaluation
script copied the fixture before running the analysis, and the copy was what
got rebuilt. The measurement was correct (the numbers described the copy
faithfully), and the measured object was wrong (the copy was not what
participants would see). The study's own participant-path guard already refused
to start a session in this state, with a one-sentence explanation of why; no
equivalent guard existed on the evaluation path, so the evaluation walked into
the state the study was protected from and reported it as the study's result.

Three instances of the same shape---a number computed correctly from an object
that was not the object it claimed to be---appeared in one day of evaluation.
The only checks that caught them were pre-declared falsification rules written
before the measurement: ``the census must rediscover the blemishes the build
log already confesses to.'' That rule found 1 of 4 known blemishes on its
first run, and the failure of the check (not the failure of the tool) was what
exposed the version mismatch. The lesson is methodological: in a system where
the analysis tool can modify the artefact being analysed, the evaluation must
verify the identity of the measured object after each run, from evidence the
tool cannot have produced. A number that passes every internal consistency
check is still wrong if the thing it describes is not the thing it claims to
describe, and the only safeguard is a rule written before the measurement that
asks ``is this the same artefact?'' rather than ``is this number consistent?''

For a CHI audience the finding connects to Sedlmair et al.'s methodological
reflection on evaluation in visualization research~\cite{sedlmair2012}: the
hardest threats to validity are not the ones that make a result look bad but
the ones that make it look good by accident, because no incentive exists to
investigate a passing check. Every one of the evaluation journal's withdrawn
results was a result that passed.

\paragraph{What the pilots could not test.}
Three things only the study with real people under time pressure can show, all
of which the pilots explicitly flagged as untestable by an AI agent:
\begin{enumerate}
\item Whether the time cap binds differently per condition (the manipulation
  check ``how much time pressure did you feel'' is the direct signal).
\item Whether carry-over between isomorphic projects is large enough to
  dominate the condition effect (order as a factor in the model).
\item Whether a person gives up --- the AI pilot participants never did, but a
  person who spends five minutes with no progress is fundamentally different
  from an agent that retries.
\end{enumerate}
\placeholder{Report: order effect size vs. condition effect size. If order
dominates, the design lesson is that counterbalancing controls for it
statistically but not experientially, and a future study should use non-isomorphic
projects despite the difficulty of equating them.}

\subsection{Scope and limitations}

Twelve participants can detect large effects and nothing else. A confidence
interval that includes zero at this sample size does not mean the effect is
absent; it means we cannot tell. We designed for this deliberately
(Section~\ref{sec:eval-study}) and report intervals rather than $p$-values so
that a reader can see the precision rather than a binary gate. The practical
consequence is that a small-to-medium benefit of the kind sgt's read layer
probably provides would appear as a wide interval tilted in the expected
direction, not as a clean separation, and the reader should treat the
qualitative pattern (which measures move, in which direction) as more
informative than any individual comparison.

Two hours with a transcript is not a week of forgetting. The study's tasks
substitute ``code you have never seen'' for ``code you wrote through an agent
and then forgot,'' and the substitution understates the difficulty of the
condition Section~\ref{sec:scenarios} describes, because a participant who just
read a transcript knows more than the developer who typed those words a week
ago. Any advantage sgt shows in this study is therefore a lower bound on the
advantage in the scenario that motivated it.

The two study repositories share the same eighteen-marker vocabulary and the
same task structure. Counterbalancing controls for which project goes first, but
a participant who solves the removal task in one repository has seen the shape
of the answer before they open the other. The order effect size
(\placeholder{report}) bounds how much of what we measured is learning the task
rather than learning the tool.

The feature labels that make the sgt condition legible depend on a language model
call. During preparation we discovered that without a credential, a rebuild
produces nine features labelled with raw symbol names (\texttt{build\_parser main
Section}) in precisely the features the tasks ask about. The shipped fixtures
pre-cache labels so that participants see domain vocabulary regardless of network
state, verified byte-identical across rebuilds with zero model calls. A view
whose legibility depends on a network call and a credential is a view that can
silently stop being the thing being evaluated; we state this as a limitation of
the design rather than of the study, because a user who installs \sgt{} without
a model key faces the same degradation and receives no warning.

The limitation that matters most is the one we cannot bound. The study measures
whether a tool helps a developer do three specific things to code they have never
seen. It cannot measure whether the same tool changes how the developer works
over weeks---whether the names enter their vocabulary, whether the attribution
shifts how they prompt the agent, whether the consequence preview changes what
they are willing to attempt. Those are the effects that would make the tool
worth its learning cost, and a two-hour session can detect none of them. What it
can detect is whether the tool's mechanisms reach the developer at all on first
contact, and the pilot showed that they can: one participant adopted the
attribution immediately and rephrased the revert as a punch list within twenty
minutes. The study cannot tell us whether that adoption persists, deepens, or
disappears once the session ends and the researcher leaves the room.
